% !TeX program = pdflatex
% !TeX TXS-program:compile = txs:///pdflatex/[-shell-escape]
% !TeX TXS-program:pdflatex = pdflatex -synctex=1 -interaction=nonstopmode --shell-escape %.tex
\section{Output And Input}\label{sec:output_and_input}

\subsection{Video - System Calls}\label{subsec:video_system_calls}
\subsection{Writing Output}\label{subsec:writing_output}

Recordando los \textit{File Descriptors(FDs)}, vemos que cada proceso empieza con alguno de estos 3 FDs:

\begin{itemize}
	\item \textbf{FD 0:} Standard Input
	\item \textbf{FD 1:} Standard Output
	\item \textbf{FD 2:} Standard Error
\end{itemize}

En el caso de la llamada al sistema para \textit{\textbf{exit}}, tenemos que
especificar que el valor de \mintinline{asm}{mov rdi, 42}. Ahora para poder
ejecutar la llamada al sistema de \textit{\textbf{write}}, tenemos que
establecer el valor de rdi en 1 o en 2.


Ahora, este no es el único parámetro que requiere la instrucción
\textit{\textbf{write}}, esta también necesita saber en que parte de la memoria
esta el texto y cuantos caracteres vamos a leer. Similar a como sería una
función en C:

\begin{codebox}[C]{Ejemplo de la función write}
	write(file_descriptor, memory address, number_of_characters_to_write)
	# Ejemplo de su uso
	write(1, memory address,10)
\end{codebox}

Ahora, por estándar, esto de aquí ya esta definido:

\begin{itemize}
	\item El primer parámetro siempre será \textit{\textbf{rdi}}.
	\item El segundo parámetro siempre será \textit{\textbf{rsi}}.
	\item El tercer parámetro siempre será \textit{\textbf{rdx}}.
\end{itemize}

\begin{codebox}[ASM]{}
	mov rax , 1
	mov rdi , 1
	mov rsi , [rsp + 16]
	mov rdx , 1
	syscall
\end{codebox}

\begin{flagbox}{Flag Capturada}
	pwn.college{MkMD83Ee-aMqL23-gYH-nLVAcvF.QXwUTN2wSMycjN4EzW}
\end{flagbox}

\subsection{Chaining Syscalls}\label{subsec:chaining_syscalls}


En el código anterior, si lo compilamos y ejecutamos, vamos a ver que el
programa se detiene, para que se no bloquee, se tiene que agregar la llamada
de salida.

\begin{codebox}[ASM]{}
	mov rax , 1
	mov rdi , 1
	mov rsi , [rsp + 16]
	mov rdx , 1
	syscall
	mov rax , 60
	mov rdi , 42
	syscall
\end{codebox}

\begin{flagbox}{Flag Capturada}
	pwn.college{k3RvbIqmOS-vpMk7rE2EgfBpya-.QXxUTN2wSMycjN4EzW}
\end{flagbox}

\subsection{Writing String}\label{subsec:writing_string}

Práctica para poder leer, más de un solo caracter, ya que en el ejemplo del
capitulo \ref{subsec:writing_output}, es para un solo digito, en este caso
solo se tendría que cambiar el valor de \textit{\textbf{rdx}}

\begin{codebox}[ASM]{}
	mov rax , 1
	mov rdi , 1
	mov rsi , [rsp + 16]
	mov rdx , 64
	syscall
	mov rax , 60
	mov rdi , 42
	syscall
\end{codebox}

\begin{flagbox}{Flag Capturada}
	pwn.college{g79tV7cz1AjIFS6t816QBs5NXqM.01NzEDMxwSMycjN4EzW}
\end{flagbox}

\subsection{Reading Data}\label{subsec:reading_data}

Ahora para poder leer, vamos a usar de ejemplo como es en el lenguaje C.

\begin{codebox}[C]{Reading syscall}
	read(0,some_address,5)
\end{codebox}

En este caso, como vemos a partir de \textit{some\_address} se va a empezar a
a rellenar, hasta leer los primeros 5 caracteres del texto.

\begin{verbatim}
       Address      | Contents
+-------------------+----------+
| some_address      | 48       |
| some_address+1    | 45       |
| some_address+2    | 4c       |
| some_address+3    | 4c       |
| some_address+4    | 4f       |
+-------------------+----------+
\end{verbatim}
Aquí un ejemplo de su como sería leer, e imprimir un texto de 128 caracteres:
\begin{codebox}[ASM]{}
	mov rax , 0
	mov rdi , 0
	mov rsi , rsp
	mov rdx , 128
	syscall

	mov rax , 1
	mov rdi , 1
	mov rsi , rsp
	mov rdx , 128
	syscall
	mov rax , 60
	mov rdi , 42
	syscall
\end{codebox}

\begin{flagbox}{Flag Capturada}
	pwn.college{8hX8fIcaF2n1GfRHFQCGNA0JmKl.0FNwUTNxwSMycjN4EzW}
\end{flagbox}

\subsection{Reading Exactly}\label{subsec:reading_exactly}

Ahora, anteriormente como hemos visto anteriormente, hemos leído y escrito
un mensaje de 128 caracteres, pero no siempre vamos a saber cuantos caracteres
puede tener el texto que vamos a leer.

La funcion \textit{\textbf{read}}, nos regresa el valor de la cantidad de
caracteres que realmente a leído, y este se almacena en \textit{\textbf{rax}}.

\begin{codebox}[ASM]{}
	mov rax , 0
	mov rdi , 0
	mov rsi , rsp
	mov rdx , 128
	syscall

	mov rdx, rax

	mov rax , 1
	mov rdi , 1
	mov rsi , rsp
	syscall
	mov rax , 60
	mov rdi , 42
	syscall
\end{codebox}

\begin{flagbox}{Flag Capturada}
	pwn.college{8tnm880ONdKT9_M4-W93McmxZWn.0VM4ATNywSMycjN4EzW}
\end{flagbox}

\subsection{Opening Files}\label{subsec:opening_files}

Para abrir un archivo, que se encuentre en algun lugar de nuestro disco, vamos
a necesitar acceso a dicho archivo, para esto necesitamos abrirlo.

\begin{codebox}[C]{Open file}
	open("/flag",0) # 0 default read-only
\end{codebox}

Ahora similar a la función \textit{\textbf{read}}, la función \textit{\textbf{open}}
nos esta devolviendo un valor en \textit{\textbf{rax}}, el cual contiene
el \textit{\textbf{fd}} del archivo que acabamos de leer, y se tiene que usar
este de aquí para poder leer el archivo.

\begin{codebox}[ASM]
	mov rax, 2
	mov rdi, [rsp+16]
	mov rsi, 0
	syscall

	mov rdi, rax
	mov rax, 0
	mov rsi, rsp
	mov rdx, 256
	syscall

	mov rdx, rax
	mov rax, 1
	mov rdi, 1
	mov rsi, rsp
	syscall

	mov rax, 60
	mov rdi, 42
	syscall
\end{codebox}

\begin{infobox}[Nota Informativa]
	En este caso, se esta pasando la dirección del archivo, como parametro,por eso es que se lee
	desde \textit{\textbf{rsp + 16}}.
\end{infobox}

\begin{flagbox}{Flag Capturada}
	pwn.college{UYeD1OtcieFHPjdDycMVcs0l_aj.0VM0czMywSMycjN4EzW}
\end{flagbox}

\subsection{Hardcoding the Filename}\label{subsec:hardcoding_the_filename}

En el caso anterio \ref{subsec:opening_files}, pasamos el valor del
archivo como parametros, pero si queremos pasar el nombre de un archivo fijo,
sería de la siguiente manera.

\begin{codebox}[ASM]
	.intel_syntax noprefix

	mov BYTE PTR [rsp], '/'
	mov BYTE PTR [rsp+1], 'f'
	mov BYTE PTR [rsp+2], 'l'
	mov BYTE PTR [rsp+3], 'a'
	mov BYTE PTR [rsp+4], 'g'
	mov BYTE PTR [rsp+5], 0

	mov rax, 2
	mov rdi, rsp
	mov rsi, 0
	syscall

	mov rdi, rax
	mov rax, 0
	mov rsi, rsp
	mov rdx, 128
	syscall

	mov rdx, rax
	mov rax, 1
	mov rdi, 1
	mov rsi, rsp
	syscall

	mov rax, 60
	mov rdi, 42
	syscall
\end{codebox}

\begin{flagbox}{Flag Capturada}
	pwn.college{ksqeBnvOJWvAvIls3V_A_gdsD9k.0lM0czMywSMycjN4EzW}
\end{flagbox}

\subsection{RIP - Relative Strings}\label{subsec:rip_relative_strings}

El problema que presentamos en la sección anterior
\ref{subsec:hardcoding_the_filename}, es que si escribimos un error al escribir
la dirección del texto, sería muy teidoso encontrar su error, o si el nombre
del archivo es muy largo.

Por lo que, en el código ensamblador, se puede contener bytes que no esten
destinados a ejecutarse, estos serían colocados después de la llamada al
sistema \textbf{\textit{exit}}.

\begin{codebox}[ASM]{Ejemplo de bits adicionales}
	_start:
	...
	mov rax, 60
	syscall
	path:
	.asciz "/flag"
\end{codebox}

En este caso, la directiva \textbf{\textit{.asciz}} va a contener los bytes de
la cadena, junto con el 0 de la termincación.

Y la etiqueta \textbf{\textit{path}} marca donde comienzan esos bytes.

Esto era factible en programas anteriores, donde los programas siempre se
cargaban en la misma dirección de memoria, por lo que se podría escribir
esto de aquí:

\begin{codebox}[ASM]{}
	_start:
	...
	mov rdi, path
	path:
	.asciz "/flag"
\end{codebox}

Esto de aquí no funciona, porque ahora los programas con compilados por un tema
de seguridad, lo que impide conocer en que parte de la memoria es que se
encuentra almacenado.

Para esto de aquí, se creo el
\textit{Instruction Pointer Relative Addressing(Dirección Relativo al Puntero de Instrucción)}
El puntero de instrucción \textbf{\textit{rip}}, va a contener la dirección
de la siguiente instrucción que se ejecutará en la CPU.

\begin{codebox}[ASM]{}
	_start:
	...
	mov rdi, [rip+path]
	...
	path:
	.asciz "/flag"
\end{codebox}

Esto último presenta el problema que el valor de \textbf{\textit{rdi}}, va a ser \textbf{\textit{"/flag"}},
no la dirección de memoria donde este esta almacenado, para esto se
tiene que utilizar el la instrucción \textbf{\textit{lea}}.

La instrucción \textbf{\textit{lea}}\textit{(Load Effective Address)}, permitiendo
que se haga referencia a que la CPU calcula todas las operacoines necesarias,
como sumar un desplamazamiento al puntero de instrucción.

\begin{codebox}[ASM]{}
	_start:
	...
	lea rdi, [rip+path]
	...
	path:
	.asciz "/flag"
\end{codebox}

\begin{codebox}[ASM]{Solución del problema}
	.intel_syntax noprefix

	mov rax, 2
	lea rdi, [rip + path]
	mov rsi, 0
	syscall

	mov rdi, rax
	mov rax, 0
	mov rsi, rsp
	mov rdx, 128
	syscall

	mov rdx, rax
	mov rax, 1
	mov rdi, 1
	mov rsi, rsp
	syscall

	mov rax, 60
	mov rdi, 42
	syscall

	path:
	.asciz "/flag"
\end{codebox}

\begin{flagbox}{Flag Capturada}
	pwn.college{c45qx3vAQJ53UI51AOM6qFuPtYE.01N3ATNywSMycjN4EzW}
\end{flagbox}
